% ----------------------
\subsection{Section 65 (30 October 1831)}
% ----------------------
	\label{DC65}
	
	% -----
	\subsubsection{Historical Background}
	% -----
		\label{DC65-HB}
		
		% --
		\paragraph{JSP}
		% --
		
			``On 30 October 1831, JS dictated a revelation later characterized as a `revelation on prayer.' William E. McLellin, who may have been present when the revelation was dictated and who made a copy of it soon after, wrote that it related to Matthew 6:10---the passage of the Lord's Prayer in which Christ petitions, `Thy kingdom come.' Although JS was working on his revision of the Bible at this time, he had already reached the last chapters of Matthew or even the first chapters of Mark. His review of Matthew 6:10 came more than six months earlier, suggesting that the revelation was not connected to the revision of that passage, although, as McLellin stated, the two certainly had a thematic tie. The revelation, which declares that the keys of God's kingdom have been returned to the earth to prepare it for the second coming of Jesus Christ, culminates with a prayer for the growth of that kingdom.
			
			``The revelation may have been dictated as part of a church service held at John and Alice (Elsa) Jacobs Johnson's home on Sunday, 30 October 1831. McLellin wrote in his journal that he preached there that day in a meeting attended by JS and several others of `the brethren \& sisters.' McLellin did not refer to this revelation in his journal, but since JS was present, this may have been the venue where the revelation was dictated.
			
			``It is unclear who served as scribe for the original manuscript. McLellin inscribed a revelation just the day before,%
				\footnote{%
					% \label{}%
					This was \hyperref[DC66]{DC 66}.
					}
			but John Whitmer, Sidney Rigdon, and Oliver Cowdery, all of whom had worked or were working as JS's scribes, were at the Johnson home that Sunday. In any case, within three weeks both Whitmer and McLellin made copies of the revelation.7 It is difficult to determine which of the two is the earliest extant version, but because Whitmer's copy (which he included in Revelation Book 1 and which is featured here) shows less punctuation, it may more closely represent the original.''%
				\footnote{%
					% \label{}%
					The JSP website doesn't have a link to the McLellin version of this revelation.
					}
		
	% -----
	\subsubsection{Versions}
	% -----
		\label{DC65-V}
		
		See the \href{https://drive.google.com/file/d/1goAvNz8jS4R8slYfMG_db55Ty2z_vmgK/view?usp=sharing}{sections comparison} document.
	
		\begin{compactdesc}
			\item[JSP] \hyperref[EMS-1831-JS-30-OCT-1831]{30 October 1831}
			\item[EMS] 1:4 (September 1832), 26
			\item[BC] ---\footnotemark
			\item[DC 1835] Section 24, 151
			\item[TS] 5:7 (1 April 1844), 482
			\item[DC 1844] Section 24, 220--221
			\item[MS] 5:11 (April 1845), 163
		\end{compactdesc}
			\footnotetext{%
				% \label{}%
				Remember that the BC ends in the middle of \hyperref[DC64]{DC 64}. See the \hyperref[DC64-V]{versions} of that revelation.
				}

	% -----
	\subsubsection{Section 65:1} % *DC65-1 
	% -----
		\label{DC65-1}

		HEARKEN, and lo, a voice as of one sent down from on high, who is mighty and powerful, whose going forth is unto the ends of the earth, yea, whose voice is unto men---Prepare ye the way of the Lord, make his paths straight.%
			\footnote{%
				% \label{}%
				Invoking language of John the Baptist here; see \hyperref[MATTHEW-3-3]{Matthew 3:3}.
				}

		% \begin{framed}
		% \scriptsize
			% \begin{compactdesc}
				% \item[JSP] 
				% % \item[BC] 
				% % \item[]
			% \end{compactdesc}
		% \end{framed}

	% -----
	\subsubsection{Section 65:2} % *DC65-2 
	% -----
		\label{DC65-2}

		The \hyperref[KEY]{keys} of the \hyperref[KINGDOM]{kingdom} of God are committed unto man on the earth, and from thence shall the gospel roll forth unto the ends of the earth, as the stone which is cut out of the mountain without hands%
			\footnote{%
				% \label{}%
				See \hyperref[DANIEL-2-45]{Daniel 2:45}.
				}
		shall roll forth, until it has filled the whole earth.

		% \begin{framed}
		% \scriptsize
			% \begin{compactdesc}
				% \item[JSP] 
				% % \item[BC] 
				% % \item[]
			% \end{compactdesc}
		% \end{framed}

	% -----
	\subsubsection{Section 65:3} % *DC65-3 
	% -----
		\label{DC65-3}

		Yea, a voice crying---Prepare ye the way of the Lord, prepare ye the supper of the Lamb,%
			\footnote{%
				% \label{}%
				See \hyperref[REVELATIONNT-19-9]{Revelation 19:9} and \hyperref[DC58-11]{DC 58:11}.
				}
		make ready for the Bridegroom.%
			\footnote{%
				% \label{}%
				See \hyperref[MATTHEW-25-10]{Matthew 25:10} and \hyperref[DC33-17]{DC 33:17}.
				}

		% \begin{framed}
		% \scriptsize
			% \begin{compactdesc}
				% \item[JSP] 
				% % \item[BC] 
				% % \item[]
			% \end{compactdesc}
		% \end{framed}

	% -----
	\subsubsection{Section 65:4} % *DC65-4 
	% -----
		\label{DC65-4}

		Pray unto the Lord, call upon his holy name, make known his wonderful works among the people.

		% \begin{framed}
		% \scriptsize
			% \begin{compactdesc}
				% \item[JSP] 
				% % \item[BC] 
				% % \item[]
			% \end{compactdesc}
		% \end{framed}

	% -----
	\subsubsection{Section 65:5} % *DC65-5 
	% -----
		\label{DC65-5}

		Call upon the Lord, that his kingdom may go forth upon the earth, that the inhabitants thereof may receive it, and be prepared for the days to come, in the which the Son of Man shall come down in heaven, clothed in the brightness of his glory, to meet the kingdom of God which is set up on the earth.

		% \begin{framed}
		% \scriptsize
			% \begin{compactdesc}
				% \item[JSP] 
				% % \item[BC] 
				% % \item[]
			% \end{compactdesc}
		% \end{framed}

	% -----
	\subsubsection{Section 65:6} % *DC65-6 
	% -----
		\label{DC65-6}

		Wherefore, may the kingdom of God go forth, that the kingdom of heaven may come, that thou, O God, mayest be glorified in heaven so on earth, that thine enemies may be subdued; for thine is the honor, power and glory,%
			\footnote{%
				% \label{}%
				For ``honor, power, and glory,'' see \hyperref[HONOR-PHRASES]{here}.
				}
		forever and ever.  Amen.

		% \begin{framed}
		% \scriptsize
			% \begin{compactdesc}
				% \item[JSP] 
				% % \item[BC] 
				% % \item[]
			% \end{compactdesc}
		% \end{framed}